\section{The Nature of Light in Quiescent Vacuum, Directional Phase Propagation, and Matter Interaction}
\label{sec:isik_yonlu_yayilim}

A central inquiry in our framework is: \textit{``If the electromagnetic vacuum is strictly quiescent in the static background and the photon carries no translational matter displacement, why does light possess a definite direction of propagation? Why does radiation from the Sun reach Earth in 8.3 minutes? What does phase transmission mean in a stationary medium?''}

\subsection{Ontological and Geometric Definition of the Photon: Phase Transmission in Stationary Vacuum}
\label{sec:foton_tanim}

\begin{definition}[The Photon: Topological Soliton Wavepacket of Quiescent Vacuum]
    \label{def:foton_tanimi_main}
    A photon is a \textbf{transverse electromagnetic phase oscillation and elastic stress wavepacket (soliton)} propagating across the stationary Euclidean vacuum continuum ($\mathcal{E}_\infty, \mathbf{v}_{\mathrm{vacuum}}=0$), triggered by quantum momentum recoil ($\mathbf{p} = \hbar \mathbf{k}$) during atomic scale transitions.
    \begin{itemize}
        \item \textbf{Stationarity of the Vacuum:} The underlying vacuum substrate does not flow or drift, and space does not expand ($\mathbf{v}_{\mathrm{space}} = 0$). Photons do not drift like inertial matter particles through space.
        \item \textbf{Speed of Light ($c$):} Because the elastic phase propagation speed of the vacuum matches the internal relativistic quantum phase ceiling of matter vortices ($v_c = 1/\sqrt{\varepsilon_0\mu_0} = c$), an observer measuring with contracting atomic rulers and clocks registers this free phase wavefront as \textbf{moving at relative velocity $c$}.
    \end{itemize}
\end{definition}

\subsection{Vacuum Quiescence and Directional Phase Waves: Conservation of Momentum}
The background space $(\mathcal{M}, \eta_{\mu\nu})$ and the electromagnetic 4-potential $A_\mu$ remain stationary, analogous to the surface of a tranquil lake. A photon is a \textbf{directional phase wavepacket} excited on this quiescent continuum.

The directional character of light originates from local momentum conservation during atomic emission:
When an excited atom (such as a hydrogen atom in the solar corona) undergoes a dipole transition, the electron accelerates in a specific non-isotropic direction. By conservation of local momentum:
\begin{equation}
    \Delta \mathbf{p}_{\mathrm{atom}} = -\hbar \mathbf{k} \implies \mathbf{p}_{\mathrm{photon}} = +\hbar \mathbf{k},
    \label{eq:momentum_korunumu}
\end{equation}
where $\mathbf{k}$ is the wave vector ($|\mathbf{k}| = \omega / v_c$). The emitted wave function is characterized by a scalar phase gradient:
\begin{equation}
    A_\mu(\mathbf{x}, t) = \epsilon_\mu e^{i \Phi(\mathbf{x}, t)}, \quad \nabla \Phi = \mathbf{k}, \quad \frac{\partial \Phi}{\partial t} = -\omega.
    \label{eq:yonlu_faz_denklemi}
\end{equation}
As emitting atoms recoil, they deposit a directional phase ripple directed toward Earth; conversely, a terrestrial laser recoils upon firing, directing a phase wavefront into deep space.

\subsection{Why the Wave Velocity Equals the Internal Phase Ceiling (\texorpdfstring{$v_c = c$}{vc})}
Regardless of the velocity with which a stone is cast into water, the propagation speed of the resulting circular ripple is dictated entirely by the medium's elastic parameters. Similarly, an atomic vortex cannot excite the vacuum continuum faster than the characteristic transmission speed $v_c = 1/\sqrt{\varepsilon_0\mu_0} = c$:
\begin{equation}
    \Box A_\mu = \left( -\frac{1}{v_c^2} \frac{\partial^2}{\partial t^2} + \nabla^2 \right) A_\mu = J_\mu.
    \label{eq:dalga_operatoru_vc}
\end{equation}
Measured against our internal atomic clock rhythm ($v_c = c$), the wavefront speed is recorded as exactly $c$.

\subsection{Mechanical Wave Analogy, Vacuum Elasticity, and the String--Knot Model}
\label{sec:tel_ve_dugum_ontolojisi}

If light is a transverse wave on a taut string, what is the string composed of?

\subsubsection{Equivalence Between Mechanical Strings and Maxwell's Equations}
The velocity of a transverse wave on a mechanical string is set by tension $T$ and linear density $\mu$:
\begin{equation}
    v_{\mathrm{string}} = \sqrt{\frac{T}{\mu}} = \sqrt{\frac{\text{Elastic Tension}}{\text{Inertia Density}}}.
    \label{eq:mekanik_tel_hizi}
\end{equation}
Maxwell's formula for the speed of light exhibits an identical mathematical structure:
\begin{equation}
    c = \frac{1}{\sqrt{\varepsilon_0 \mu_0}} = \sqrt{\frac{1/\varepsilon_0}{\mu_0}} = \sqrt{\frac{\text{Vacuum Electric Bulk Modulus}}{\text{Vacuum Magnetic Inertia Density}}}.
    \label{eq:vakum_tel_ozdeslik}
\end{equation}
\begin{itemize}
    \item \textbf{Electric Modulus ($1/\varepsilon_0$):} Quantifies the vacuum's elastic resistance against electric polarization and displacement.
    \item \textbf{Magnetic Inertia Density ($\mu_0$):} Represents the vacuum's self-inductive opposition to magnetic field changes.
    \item \textbf{Characteristic Impedance ($Z_0 \approx 377\,\Omega$):} Vacuum is not an absolute void; it possesses a measurable electromagnetic impedance:
          \begin{equation}
              Z_0 \equiv \sqrt{\frac{\mu_0}{\varepsilon_0}} \approx 376.73\,\Omega.
              \label{eq:vakum_empedansi_tanim}
          \end{equation}
\end{itemize}

\subsubsection{String and Knot Ontology: Unified Origin of Matter and Photons}
\begin{itemize}
    \item \textbf{The String (Vacuum Continuum):} The electromagnetic 4-potential field $A_\mu$ defined over static Euclidean space $\mathcal{E}_\infty$.
    \item \textbf{The Photon (Free Wave):} A transverse phase ripple traveling along the unknotted string at speed $c$ ($m_\gamma = 0$).
    \item \textbf{Matter (Topological Knot):} A localized, closed chiral vortex knot in the continuum lines (soliton). Knots localize rotational phase, conferring rest mass ($E=mc^2$). Moving a knot through the continuum encounters inertial drag; hence $v_{\mathrm{trans}} < c$.
\end{itemize}

\subsection{Sun-to-Earth Propagation: The 8.3-Minute Signal Delay}
\label{sec:kavramsal_netlik_main}

The 8.3-minute transit time follows from a deterministic four-stage causal chain:
\begin{enumerate}
    \item \textbf{Emission on the Sun:} An excited hydrogen atom undergoes an electronic transition, recoiling and imparting a directional phase stress ($\mathbf{k}$) onto the quiescent vacuum.
    \item \textbf{Phase Transmission Through Space:} Space between the Sun and Earth does not move ($\mathbf{v}_{\mathrm{space}} = 0$). Adjacent quantum cells of the quiescent vacuum transfer the phase perturbation, propagating at $c = 1/\sqrt{\varepsilon_0\mu_0} \approx 300{,}000\text{ km/s}$.
    \item \textbf{8.3-Minute Delay (Static Euclidean Distance):} The static Euclidean distance is $d = 1.496 \times 10^{11}\text{ m}$. Traversing this interval requires:
          \begin{equation}
              \Delta t = \frac{d}{c} = \frac{1.496 \times 10^{11}\text{ m}}{2.998 \times 10^8\text{ m/s}} \approx 499\text{ seconds} \approx 8.3\text{ minutes}.
          \end{equation}
          Light-cone causality is preserved.
    \item \textbf{Absorption on Earth (Phase-Pinning):} Atoms in a leaf or retina absorb the resonant wavefront ($\hbar\omega = \Delta E, v_c = c$), collapsing the wave via phase-pinning into a localized excitation.
\end{enumerate}

\subsection{Velocity Hierarchy: Matter ($v < c$) vs. Light ($v = c$)}
\label{sec:hiz_hiyerarsisi_cozum}

\begin{itemize}
    \item \textbf{Massive Matter (Scalor Knots -- $0 \le v < c$):} A closed vortex knot displaced across the quiescent vacuum lattice experiences inertia; hence $v < c$.
    \item \textbf{Photons (Free Phase Waves -- $v_{\mathrm{phase}} = c$):} Being an elastic deformation of the medium itself, the wave propagates at the medium's characteristic speed $c$.
    \item \textbf{Coordinate Separation Speed ($v_{\mathrm{sep}} \le 2c$):} Two diametrically opposed wavefronts separate at Euclidean rate $2c$ across the static background without violating relativistic causality.
\end{itemize}

\begin{table}[h!]
    \centering
    \small
    \resizebox{\textwidth}{!}{
    \begin{tabular}{@{}p{3.0cm}p{4.5cm}p{4.0cm}p{3.8cm}@{}}
        \toprule
        \textbf{Regime}                                    & \textbf{Structure}                                                                              & \textbf{Motion}                                                                       & \textbf{Velocity Limit}                                                                                            \\ \midrule
        \textbf{Massive Matter} \newline (Electron, Proton) & Closed topological vortex knot (Scalor knot).                                                   & Displaces across the stationary vacuum grid.                                          & Possesses inertia and rest mass; strictly $v_{\mathrm{trans}} < c$.                                                \\ \addlinespace
        \textbf{Photon} \newline (Electromagnetic Wave)    & Free elastic vibration across neighboring vacuum cells; NOT a localized knot.                   & No translational matter displacement ($v_{\mathrm{trans}} = 0$); phase propagates.   & Recorded as $v_{\mathrm{phase}} = c$ relative to contracting observers.                                            \\ \addlinespace
        \textbf{Two Opposed Photons}                       & Two wavefronts traveling at $180^\circ$ separation.                                             & Rate of Euclidean distance opening.                                                   & Coordinate separation velocity $v_{\mathrm{sep}} = 2c$ (no local particle exceeds $c$).                          \\ \bottomrule
    \end{tabular}
    }
    \caption{Structural, kinematic, and velocity comparison between matter, photons, and opposing wavefronts.}
    \label{tab:rejim_karsilastirma}
\end{table}

\subsection{Why Photons Do Not Contract: Precondition for Redshift}
\label{sec:foton_neden_buzulmez}

\begin{theorem}[Precondition for Cosmological Redshift]
    \label{thm:kirmiziya_kayma_varlik_sarti}
    If photons contracted synchronously with bound matter ($\lambda_{\mathrm{photon}}(t) \propto \lambda(t)$), the dimensionless wavelength ratio measured by contracting atomic rulers ($\lambda_{\mathrm{atom}}(t) \propto \lambda(t)$) would be:
    \begin{equation}
        \frac{\lambda_{\mathrm{photon}}(t)}{\lambda_{\mathrm{atom}}(t)} = \frac{\lambda_0 e^{-H_0 t}}{r_0 e^{-H_0 t}} = \frac{\lambda_0}{r_0} = \text{const} \implies z \equiv 0.
    \end{equation}
    Cosmological redshift would be \textbf{strictly unobservable!} The necessary and sufficient condition for observing cosmic redshift ($1+z = \lambda_{\mathrm{obs}}/\lambda_{\mathrm{lab}}$) is that the photon preserves its absolute Euclidean wavelength in the stationary vacuum ($\lambda_{\mathrm{photon}} = \text{const}$), while atomic rulers continuously contract ($\lambda_{\mathrm{atom}}(t) \to 0$).
\end{theorem}

\begin{proof}
    Massless particles ($m_\gamma = 0$) do not undergo conformal mass scaling ($0 \times e^{H_0 t} = 0$). Photons deposited in the quiescent vacuum preserve their invariant Euclidean wavelength ($\lambda_{\text{real}}$). Only the measuring atomic rulers contract.
\end{proof}

\subsection{Wavelength Scale of Cosmic Radiation}
\label{sec:fotonlar_neden_devasa_degil}

The wavelength of an emitted photon is governed by the dimensions of the emitting atom:
\begin{enumerate}
    \item \textbf{Atomic Size Sets Emitted Wavelength:} The energy difference between atomic levels scales as $\Delta E(t) \propto e^2 / r_{\mathrm{atom}}(t)$, yielding $\lambda_{\mathrm{photon, emission}}(t) \propto r_{\mathrm{atom}}(t)$.
    \item \textbf{Contemporary Photons:} Emitted by today's shrunken atoms ($r_{\mathrm{atom}} \sim 10^{-10}\text{ m}$), falling into the optical band ($\lambda \sim 400\text{--}700\text{ nm}$).
    \item \textbf{Ancient Photons:} Emitted by larger atoms in the past, appearing redshifted when measured with present-day smaller rulers.
    \item \textbf{The Oldest Photons (CMB, $z \sim 1100$):} Emitted when atoms were $1100$ times larger; original ultraviolet/visible radiation ($\sim 1\text{ }\mu\mathrm{m}$) is measured today at \textbf{$\lambda \sim 1\text{ mm}$ ($2.73\text{ K}$)}.
\end{enumerate}
